\documentclass[conference]{IEEEtran}
\usepackage{graphicx}
\usepackage{amsmath}
\usepackage{array}
\usepackage{url}
\usepackage{xcolor}

% IEEEtran asks for Times (ptm), which has no TU-encoded shapes under XeTeX, so
% every font silently falls back to Latin Modern and the bold shape is dropped.
% These three lines are what keep bold, small caps and monospace working.
\renewcommand{\rmdefault}{lmr}
\renewcommand{\sfdefault}{lmss}
\renewcommand{\ttdefault}{lmtt}

\begin{document}

\title{FSD-Home: Full Self-Driving Home Orchestration\\ with Risk-Graded Autonomous Action}

\author{
\IEEEauthorblockN{Ha Yun Kim, SooYoung Song, Ruitan Sun, Tang Bin}
\IEEEauthorblockA{Department of Information Systems\\
Hanyang University, Seoul, Korea\\
hykim1215@hanyang.ac.kr, sys1109@hanyang.ac.kr,\\
sunruitan2024@126.com, tbrobo99@outlook.com}
}
\maketitle

\begin{abstract}
Most smart home setups force users to open separate apps for every appliance or
react to push notifications after a problem has already occurred. We built
FSD-Home (Full Self-Driving Home) to make home automation proactive, inclusive,
and hassle-free for all household members [1], [2]. Our system continuously
monitors simulated IoT telemetry---like room temperature, laundry humidity,
compressor loads, and daily schedules---to catch risks early, such as mold
growing in a closed washer or a kitchen overheating during cooking [3], [4]. To
keep the system safe and predictable across the household, we use a Level 1 /
Level 2 autonomy framework: low-risk adjustments (micro-ventilation,
pre-cooling) run automatically in the background with a brief notice, while
high-stakes actions (delaying a high-power dryer or preheating a cooktop)
require clear approval on screen or by voice [5]. By decoupling a fast local
rule check (1.2 ms) from LLM reasoning (2--3 s), FSD-Home runs safely and
reliably without stalling the home state.
\end{abstract}

\begin{IEEEkeywords}
Smart Home Orchestration, FSD-Home, Full Self Driving Home, Risk-Graded
Autonomy, Level 1 and Level 2 Control, Inclusive Usability, Energy
Load-Shifting, LG ThinQ Ecosystem.
\end{IEEEkeywords}

\begin{center}
\centering
\footnotesize
{\scshape\bfseries TABLE I\par}
{\scshape\bfseries ROLE ASSIGNMENTS\par}
\vspace{2pt}
\begin{tabular}{@{}>{\raggedright\arraybackslash}p{\dimexpr 0.235\columnwidth-3.4pt\relax}
>{\raggedright\arraybackslash}p{\dimexpr 0.175\columnwidth-4.8pt\relax}
>{\raggedright\arraybackslash}p{\dimexpr 0.53\columnwidth-14.4pt\relax}@{}}
\hline
\textbf{Role} & \textbf{Assigned Member} & \textbf{Task Description \& Responsibilities} \\
\hline
User & Ruitan Sun & Identify user needs, provide feedback, and help ensure the service meets user expectations. \\
Customer & SooYoung Song & Organize customer and household requirements and defines safety boundaries and sets Level 2 confirmation policy rules. \\
Software Developer & Tang Bin & Builds the core AI decision engine (Python), feeds data, constructs integration tools, and generates simulated IoT sensor data. \\
Development Manager & Ha Yun Kim & Leads system architecture design, manages UI design, maintains project documentation, Git workflows, and tracks milestone progress. \\
\hline
\end{tabular}
\end{center}

\section{Introduction}

\subsection{Motivation}
Smart home appliances are ubiquitous today, but having more features does not
automatically ensure a seamless experience for every member of a household [9].
Complex menus, fragmented mobile applications, and lack of cross device
coordination create unnecessary cognitive burden [10]. Over 67\% of surveyed
consumers report experiencing friction with hidden menus and unclear controls on
modern appliances [11]. Furthermore, routine household oversights cause major
safety and hygiene issues: 35.3\% of residential fires start from unattended
cooking [12], [13], and wet laundry left in closed washers quickly leads to mold
and bacterial growth [3], [14].

\subsection{Problem Statement}
Current smart home apps like LG ThinQ work in reactive silos [15]. They wait for
manual button taps or send passive alerts (such as ``Cycle Finished''), leaving
household members to figure out what to do [6]. More importantly, appliances do
not talk to each other about room heat or power load. When someone bakes in a hot
oven while running a dryer, room temperatures surge, forcing the refrigerator and
AC to work harder and draw peak electricity without any coordinated cooling [4].
FSD-Home solves this by combining continuous sensor tracking with automated
cross-appliance control.

\subsection{Related Work \& Commercial Baseline}
Current commercial tools offer great single-device features: LG AI DD\texttrademark{}
senses fabric softness via motor torque; LG ezDispense\texttrademark{} handles
detergent dosing; and LG ThinQ ON / ThinQ Claw provides conversational voice hubs
[1], [15]. Samsung's Bespoke AI sends lifestyle pattern notifications [2], [16].
But none of them coordinate temperature and energy across multiple devices [5].
FSD-Home introduces three main contributions: (1) explicit risk scores for mold
(MRI) and kitchen heat (TOI), (2) a Tesla-inspired Level 1 / Level 2 autonomy
model with safe standby fallback, and (3) an inclusive overlay designed around
WCAG 2.2 contrast and touch targets [17].

\subsection{Academic Foundations}
We drew key design principles from recent smart home research. Hewitt (2026)
showed that conversational agents need strict deterministic gates to verify AI
outputs before triggering hardware [8]. Yuan et al.\ (2024; R-Judge) proved that
giving an LLM direct execution authority causes safety risks (like accidentally
shutting off power) [5]. Zhang et al.\ (2026) pointed out that safety and utility
metrics must be evaluated separately [7]. Following these insights, FSD-Home uses
AI only to interpret intent and write natural explanations, while our local
Python rule engine retains 100\% execution authority [18].

\section{System Requirements}

\subsection{Functional Requirements}
\begin{itemize}
  \item \textbf{FR-1: Real-Time Telemetry Perception.} The system continuously
  ingests simulated 1 Hz IoT sensor streams (temperatures, drum humidity decay,
  power draw, door contact states, calendar schedules) [19].
  \item \textbf{FR-2: Predictive Risk Calculation.} The system computes two
  real-time risk metrics in \texttt{app/risk.py}:
  \begin{enumerate}
    \item \textbf{Mold Risk Index (MRI):} tracks trapped moisture in sealed drum
    appliances using a 30-minute regression extrapolation [3]:
    \begin{equation*}
    \begin{split}
      \mathrm{MRI} = \mathrm{clamp}\big(&(\text{humidity}-70.0)\times 3.2\\
        &+ (\text{sealed\_hours}\times 6.0)\\
        &+ \max(0.0, \text{temp}-24.0)\times 1.5,\\
        &0,\ 100\big)
    \end{split}
    \end{equation*}
    \item \textbf{Thermal Overload Index (TOI):} tracks heat buildup and
    compressor strain across kitchen appliances [4]:
    \begin{equation*}
    \begin{split}
      \mathrm{TOI} = \mathrm{clamp}\big(&\text{heat\_w}/25.0\\
        &+ \max(0.0, \text{temp}-27.0)\times 6.0\\
        &+ \max(0.0, \text{duty}-0.6)\times 60.0\\
        &+ \text{cooking\_boost},\ 0,\ 100\big)
    \end{split}
    \end{equation*}
  \end{enumerate}
  When either score hits or exceeds 50.0, the system evaluates action candidates
  [5].
  \item \textbf{FR-3: Level 1 \& Level 2 Action Framework.}
  \begin{enumerate}
    \item \textbf{Level 1 (Low-Risk Autonomous Actions):} low power background
    tweaks (drum micro-ventilation, air purifier redirection, range hood
    extraction, pre-cooling) run automatically and log a plain-language summary [8].
    \item \textbf{Level 2 (High-Risk Approval Actions):} high-power operations or
    schedule changes (shifting heavy dryer runs to off-peak hours, cooktop
    preheating) strictly require user approval first [7].
  \end{enumerate}
  \item \textbf{FR-4: Inclusive Conversational Interface.} Prompts are spoken via
  Text-to-Speech (TTS) and displayed on high contrast overlays with large touch
  targets ($\geq 24 \times 24$ px) accessible for all users [17], [20].
  \item \textbf{FR-5: Fail-Safe Timeout.} If a Level 2 prompt gets no response
  within 5 minutes (300 s), the request safely expires and leaves device
  schedules unchanged [5].
\end{itemize}

\subsection{Non-Functional Requirements}
\begin{itemize}
  \item \textbf{NFR-1: Processing Latency.} The local risk check
  (\texttt{risk.py}) executes in $<1.2$ ms per cycle. Asynchronous LLM reasoning
  finishes in 2--3 seconds upon a risk trigger, controlled by a 180 s cooldown
  window [8].
  \item \textbf{NFR-2: Local Simulation Containment.} All sensor streams and
  device states run locally using generated JSONL/CSV log files [21].
  \item \textbf{NFR-3: Strict Whitelist Safety.} The decision engine enforces a
  strict 7-action whitelist, guaranteeing zero unapproved Level 2 execution [7].
\end{itemize}

\begin{thebibliography}{99}
\footnotesize
\bibitem{r1} LG Electronics, ``LG Electronics Unveils the Next Evolution of AI Home at IFA 2026,'' Global Newsroom, Aug. 2026.
\bibitem{r2} Samsung Electronics, ``Home Companion: Comprehension of Samsung Bespoke AI 2026,'' Jan. 2026.
\bibitem{r3} U.S. Environmental Protection Agency (EPA), ``A Brief Guide to Mold, Moisture and Your Home,'' EPA 402-K-02-003, 2022.
\bibitem{r4} ASHRAE, ``Standard 55: Thermal Environmental Conditions for Human Occupancy,'' ASHRAE, 2023.
\bibitem{r5} T. Yuan et al., ``R-Judge: Benchmarking Safety Risk Awareness for LLM Agents,'' arXiv:2401.10019, 2024.
\bibitem{r6} A. Pradhan, A. Lazar, and L. Findlater, ``Use of Intelligent Voice Assistants by Older Adults,'' ACM TOCHI, vol. 27, no. 4, 2020.
\bibitem{r7} H. Zhang et al., ``PromptShield Home: Ambient Multimodal Prompt Injection Defense,'' arXiv:2608.05495, 2026.
\bibitem{r8} M. Hewitt, ``Speaking to the Edge: Conversational Control of Smart Home IoT,'' Ph.D. thesis, Univ. of Sheffield, 2026.
\bibitem{r9} Ministry of Science and ICT, ``2025 Survey on the Digital Divide,'' KDI Economic Info Center, Mar. 2026.
\bibitem{r10} Money Today, ``2025 Digital Divide Survey: Competency Level Analysis,'' Mar. 2026.
\bibitem{r11} Maeil Business Newspaper, ``Home Appliance Usability Survey for Consumers 50+,'' 2026.
\bibitem{r12} Seoul Metropolitan Fire Headquarters, ``2019 Residential Fire Statistics Analysis,'' 2020.
\bibitem{r13} National Fire Agency, ``Chuseok Holiday Fire Safety Bulletin,'' Seoul Shinmun, 2025.
\bibitem{r14} World Health Organization (WHO), ``WHO Guidelines for Indoor Air Quality: Dampness and Mould,'' 2009.
\bibitem{r15} LG Electronics, ``LG ThinQ Platform Specifications,'' Official Whitepaper, 2026.
\bibitem{r16} Samsung India, ``2026 Range of Bespoke AI Appliances,'' Press Release, 2026.
\bibitem{r17} W3C, ``Web Content Accessibility Guidelines (WCAG) 2.2,'' W3C Recommendation, 2023.
\bibitem{r18} AdaHome, ``An Adaptive Smart Home Assistant using Local Language Models,'' arXiv:2607.18034, 2026.
\bibitem{r19} Korea Meteorological Administration, ``Open Data Portal APIs,'' 2026.
\bibitem{r20} C. Murad et al., ``Design guidelines for hands-free speech interaction,'' in Proc. MobileHCI '18, 2018.
\bibitem{r21} NIST, ``IoT Device Cybersecurity Guidance,'' NIST Special Publication 800-213, 2021.
\end{thebibliography}

\end{document}
